\begin{thebibliography}{99}
\footnotesize\itemsep=0pt

\bibitem{AM}
Adamovi\'{c} D., Milas A., On some vertex algebras related to
 {$V_{-1}(\mathfrak{sl}(n))$} and their characters,
 \href{https://doi.org/10.1007/s00031-020-09617-w}{\textit{Transform. Groups}}
 \textbf{26} (2021), 1--30,
 \href{https://arxiv.org/abs/1805.09771}{arXiv:1805.09771}.

\bibitem{Andrews0}
Andrews G.E., The theory of partitions, \textit{Encyclopedia Math. Appl.},
 Vol.~2, Addison-Wesley Publishing Co., Reading, Mass., 1976.

\bibitem{Andrews}
Andrews G.E., Stacked lattice boxes,
 \href{https://doi.org/10.1007/BF01608779}{\textit{Ann. Comb.}} \textbf{3}
 (1999), 115--130.

\bibitem{AR}
Andrews G.E., Rose S.C.F., Mac{M}ahon's sum-of-divisors functions, {C}hebyshev
 polynomials, and quasi-modular forms,
 \href{https://doi.org/10.1515/crelle.2011.179}{\textit{J.~Reine Angew.
 Math.}} \textbf{676} (2013), 97--103,
 \href{https://arxiv.org/abs/1010.5769}{arXiv:1010.5769}.

\bibitem{Arakawa}
Arakawa T., Chiral algebras of class $\mathcal{S}$ and {M}oore--{T}achikawa
 symplectic varieties,
 \href{https://arxiv.org/abs/1811.01577}{arXiv:1811.01577}.

\bibitem{AK}
Arakawa T., Kawasetsu K., Quasi-lisse vertex algebras and modular linear
 differential equations, in Lie Groups, Geometry, and Representation Theory,
 \textit{Progr. Math.}, Vol. 326,
 \href{https://doi.org/10.1007/978-3-030-02191-7_2}{Birkh\"auser}, Cham, 2018,
 41--57, \href{https://arxiv.org/abs/1610.05865}{arXiv:1610.05865}.

\bibitem{BK2}
Bachmann H., K\"uhn U., A short note on a conjecture of {O}kounkov about a
 $q$-analogue of multiple zeta values,
 \href{https://arxiv.org/abs/1407.6796}{arXiv:1407.6796}.

\bibitem{BK1}
Bachmann H., K\"uhn U., The algebra of generating functions for multiple
 divisor sums and applications to multiple zeta values,
 \href{https://doi.org/10.1007/s11139-015-9707-7}{\textit{Ramanujan~J.}}
 \textbf{40} (2016), 605--648,
 \href{https://arxiv.org/abs/1309.3920}{arXiv:1309.3920}.

\bibitem{BK3}
Bachmann H., K\"uhn U., A dimension conjecture for $q$-analogues of multiple
 zeta values, in Periods in Quantum Field Theory and Arithmetic,
 \textit{Springer Proc. Math. Stat.}, Vol.~314,
 \href{https://doi.org/10.1007/978-3-030-37031-2_9}{Springer}, Cham, 2020,
 237--258, \href{https://arxiv.org/abs/1708.07464}{arXiv:1708.07464}.

\bibitem{BR}
Beem C., Lemos M., Liendo P., Peelaers W., Rastelli L., van Rees B.C., Infinite
 chiral symmetry in four dimensions,
 \href{https://doi.org/10.1007/s00220-014-2272-x}{\textit{Comm. Math. Phys.}}
 \textbf{336} (2015), 1359--1433,
 \href{https://arxiv.org/abs/1312.5344}{arXiv:1312.5344}.

\bibitem{Beem}
Beem C., Rastelli L., Vertex operator algebras, {H}iggs branches, and modular
 differential equations,
 \href{https://doi.org/10.1007/jhep08(2018)114}{\textit{J.~High Energy Phys.}}
 \textbf{2018} (2018), no.~8, 114, 72~pages,
 \href{https://arxiv.org/abs/1707.07679}{arXiv:1707.07679}.

\bibitem{Beem-new}
Beem C., Singh P., Razamat S.S., Schur indices of class {$\mathcal{S}$} and
 quasimodular forms,
 \href{https://doi.org/10.1103/physrevd.105.085009}{\textit{Phys. Rev.~D}}
 \textbf{105} (2022), 085009, 15~pages,
 \href{https://arxiv.org/abs/2112.10715}{arXiv:2112.10715}.

\bibitem{Bradley}
Bradley D.M., Multiple {$q$}-zeta values,
 \href{https://doi.org/10.1016/j.jalgebra.2004.09.017}{\textit{J.~Algebra}}
 \textbf{283} (2005), 752--798,
 \href{https://arxiv.org/abs/math/0402093}{arXiv:math/0402093}.

\bibitem{BFM}
Bringmann K., Folsom A., Mahlburg K., Quasimodular forms and
 {$s\ell(m|m)^\wedge$} characters,
 \href{https://doi.org/10.1007/s11139-014-9621-4}{\textit{Ramanujan~J.}}
 \textbf{36} (2015), 103--116,
 \href{https://arxiv.org/abs/1302.4040}{arXiv:1302.4040}.

\bibitem{BJM}
Bringmann K., Jennings-Shaffer C., Milas A., Graph schemes, graph series, and
 modularity,
 \href{https://doi.org/10.1016/j.jcta.2023.105749}{\textit{J.~Combin. Theory
 Ser.~A}} \textbf{197} (2023), 105749, 37~pages,
 \href{https://arxiv.org/abs/2105.05660}{arXiv:2105.05660}.

\bibitem{BVM}
Bringmann K., Milas A., van Ittersum J.W.M., Quasimodularity  of $q$-series associated to root lattices, {i}n preparation.

\bibitem{CLM}
Calinescu C., Lepowsky J., Milas A., Vertex-algebraic structure of the
 principal subspaces of level one modules for the untwisted affine {L}ie
 algebras of types {$A$, $D$, $E$},
 \href{https://doi.org/10.1016/j.jalgebra.2009.09.029}{\textit{J.~Algebra}}
 \textbf{323} (2010), 167--192,
 \href{https://arxiv.org/abs/0908.4054}{arXiv:0908.4054}.

\bibitem{Carlsson}
Carlsson E., Vertex operators, {G}rassmannians, and {H}ilbert schemes,
 \href{https://doi.org/10.1007/s00220-010-1123-7}{\textit{Comm. Math. Phys.}}
 \textbf{300} (2010), 599--613,
 \href{https://arxiv.org/abs/0910.5528}{arXiv:0910.5528}.

\bibitem{DMN}
Dong C., Mason G., Nagatomo K., Quasi-modular forms and trace functions
 associated to free boson and lattice vertex operator algebras,
 \href{https://doi.org/10.1155/S1073792801000204}{\textit{Int. Math. Res.
 Not.}} \textbf{2001} (2001), 409--427,
 \href{https://arxiv.org/abs/math/0004132}{arXiv:math/0004132}.

\bibitem{DFR}
Dotsenko V., Feigin E., Reineke M., Koszul algebras and {D}onaldson--{T}homas
 invariants, \href{https://doi.org/10.1007/s11005-022-01604-4}{\textit{Lett.
 Math. Phys.}} \textbf{112} (2022), 106, 39~pages,
 \href{https://arxiv.org/abs/2111.07588}{arXiv:2111.07588}.

\bibitem{DM}
Dotsenko V., Mozgovoy S., D{T} invariants from vertex algebras,
 \href{https://doi.org/10.1017/S1474748024000288}{\textit{J.~Inst. Math.
 Jussieu}} \textbf{24} (2025), 291--339,
 \href{https://arxiv.org/abs/2108.10338}{arXiv:2108.10338}.

\bibitem{Efimov}
Efimov A.I., Cohomological {H}all algebra of a symmetric quiver,
 \href{https://doi.org/10.1112/S0010437X12000152}{\textit{Compos. Math.}}
 \textbf{148} (2012), 1133--1146,
 \href{https://arxiv.org/abs/1103.2736}{arXiv:1103.2736}.

\bibitem{Ekholm}
Ekholm T., Gruen A., Gukov S., Kucharski P., Park S., Sto\v{s}i\'{c} M.,
 Su{\l}kowski P., Branches, quivers, and ideals for knot complements,
 \href{https://doi.org/10.1016/j.geomphys.2022.104520}{\textit{J.~Geom.
 Phys.}} \textbf{177} (2022), 104520, 75~pages,
 \href{https://arxiv.org/abs/2110.13768}{arXiv:2110.13768}.

\bibitem{FS}
Feigin B., Stoyanovsky A.V., Quasi-particles models for the representations of
 {L}ie algebras and geometry of flag manifold,
 \href{https://arxiv.org/abs/hep-th/9308079}{arXiv:hep-th/9308079}.

\bibitem{Goncharov}
Goncharov A.B., Multiple polylogarithms, cyclotomy and modular complexes,
 \href{https://doi.org/10.4310/MRL.1998.v5.n4.a7}{\textit{Math. Res. Lett.}}
 \textbf{5} (1998), 497--516.

\bibitem{GM}
Goujard E., M\"{o}ller M., Counting {F}eynman-like graphs: {Q}uasimodularity
 and {S}iegel--{V}eech weight,
 \href{https://doi.org/10.4171/jems/924}{\textit{J.~Eur. Math. Soc. (JEMS)}}
 \textbf{22} (2020), 365--412,
 \href{https://arxiv.org/abs/1609.01658}{arXiv:1609.01658}.

\bibitem{CM}
Herzog J., Hibi T., Distributive lattices, bipartite graphs and {A}lexander
 duality,
 \href{https://doi.org/10.1007/s10801-005-4528-1}{\textit{J.~Algebraic
 Combin.}} \textbf{22} (2005), 289--302,
 \href{https://arxiv.org/abs/math/0307235}{arXiv:math/0307235}.

\bibitem{Hoffman}
Hoffman M.E., Algebraic aspects of multiple zeta values, in Zeta Functions,
 Topology and Quantum Physics, \textit{Dev. Math.}, Vol.~14,
 \href{https://doi.org/10.1007/0-387-24981-8_4}{Springer}, New York, 2005,
 51--73, \href{https://arxiv.org/abs/math/0309425}{arXiv:math/0309425}.

\bibitem{Zagier1}
Ihara K., Kaneko M., Zagier D., Derivation and double shuffle relations for
 multiple zeta values,
 \href{https://doi.org/10.1112/S0010437X0500182X}{\textit{Compos. Math.}}
 \textbf{142} (2006), 307--338.

\bibitem{JM1}
Jennings-Shaffer C., Milas A., On {$q$}-series identities for false theta
 series, \href{https://doi.org/10.1016/j.aim.2020.107411}{\textit{Adv. Math.}}
 \textbf{375} (2020), 107411, 22~pages,
 \href{https://arxiv.org/abs/2001.11368}{arXiv:2001.11368}.

\bibitem{JM2}
Jennings-Shaffer C., Milas A., Further {$q$}-series identities and conjectures
 relating false theta functions and characters, in Lie groups, Number Theory,
 and Vertex Algebras, \textit{Contemp. Math.}, Vol.~768,
 \href{https://doi.org/10.1090/conm/768/15467}{American Mathematical Society},
 Providence, RI, 2021, 253--269,
 \href{https://arxiv.org/abs/2005.13620}{arXiv:2005.13620}.

\bibitem{Kawasetsu}
Kawasetsu K., The free generalized vertex algebras and generalized principal
 subspaces,
 \href{https://doi.org/10.1016/j.jalgebra.2015.07.014}{\textit{J.~Algebra}}
 \textbf{444} (2015), 20--51,
 \href{https://arxiv.org/abs/1502.05276}{arXiv:1502.05276}.

\bibitem{Quivers}
Kucharski P., Reineke M., Sto\v{s}i\'{c} M., Su{\l}kowski P., Knots-quivers
 correspondence,
 \href{https://doi.org/10.4310/atmp.2019.v23.n7.a4}{\textit{Adv. Theor. Math.
 Phys.}} \textbf{23} (2019), 1849--1902,
 \href{https://arxiv.org/abs/1707.04017}{arXiv:1707.04017}.

\bibitem{Lang}
Lang S., Elliptic functions, 2nd ed., \textit{Grad. Texts in Math.}, Vol.~112,
 \href{https://doi.org/10.1007/978-1-4612-4752-4}{Springer}, New York, 1987.

\bibitem{Hao}
Li H., Some remarks on associated varieties of vertex operator superalgebras,
 \href{https://doi.org/10.1007/s40879-021-00477-6}{\textit{Eur.~J. Math.}}
 \textbf{7} (2021), 1689--1728,
 \href{https://arxiv.org/abs/2007.04522}{arXiv:2007.04522}.

\bibitem{HM}
Li H., Milas A., Jet schemes, quantum dilogarithm and {F}eigin--{S}toyanovsky's
 principal subspaces,
 \href{https://doi.org/10.1016/j.jalgebra.2023.10.033}{\textit{J.~Algebra}}
 \textbf{640} (2024), 21--58,
 \href{https://arxiv.org/abs/2010.02143}{arXiv:2010.02143}.

\bibitem{Libgober}
Libgober A., Elliptic genera, real algebraic varieties and quasi-{J}acobi
 forms, in Topology of Stratified Spaces, \textit{Math. Sci. Res. Inst.
 Publ.}, Vol.~58, Cambridge University Press, Cambridge, 2011, 95--120,
 \href{https://arxiv.org/abs/0904.1026}{arXiv:0904.1026}.

\bibitem{M2}
Milas A., Formal differential operators, vertex operator algebras and
 zeta-values,~{II},
 \href{https://doi.org/10.1016/S0022-4049(03)00036-7}{\textit{J.~Pure Appl.
 Algebra}} \textbf{183} (2003), 191--244,
 \href{https://arxiv.org/abs/math/0303154}{arXiv:math/0303154}.

\bibitem{M1}
Milas A., On certain automorphic forms associated to rational vertex operator
 algebras, in Moonshine: The First Quarter Century and Beyond, \textit{London
 Math. Soc. Lecture Note Ser.}, Vol.~372,
 \href{https://doi.org/10.1017/CBO9780511730054.016}{Cambridge University
 Press}, Cambridge, 2010, 330--357.

\bibitem{MP}
Milas A., Penn M., Lattice vertex algebras and combinatorial bases: general
 case and {$\mathcal W$}-algebras, \textit{New York J. Math.} \textbf{18}
 (2012), 621--650.

\bibitem{Ob-Serre}
Oberdieck G., A {S}erre derivative for even weight {J}acobi forms,
 \href{https://arxiv.org/abs/1209.5628}{arXiv:1209.5628}.

\bibitem{Ob}
Oberdieck G., Gromov--{W}itten invariants of the {H}ilbert schemes of points of
 a {K}3 surface, \href{https://doi.org/10.2140/gt.2018.22.323}{\textit{Geom.
 Topol.}} \textbf{22} (2018), 323--437,
 \href{https://arxiv.org/abs/1406.1139}{arXiv:1406.1139}.

\bibitem{Ohno}
Ohno Y., Okuda J., Zudilin W., Cyclic {$q$}-{MZSV} sum,
 \href{https://doi.org/10.1016/j.jnt.2011.08.001}{\textit{J.~Number Theory}}
 \textbf{132} (2012), 144--155.

\bibitem{Zagier2}
Ohno Y., Zagier D., Multiple zeta values of fixed weight, depth, and height,
 \href{https://doi.org/10.1016/S0019-3577(01)80037-9}{\textit{Indag.
 Math.~(N.S.)}} \textbf{12} (2001), 483--487.

\bibitem{Okounkov}
Okunkov A., Hilbert schemes and multiple {$q$}-zeta values,
 \href{https://doi.org/10.1007/s10688-014-0054-z}{\textit{Funct. Anal. Appl.}}
 \textbf{48} (2014), 138--144,
 \href{https://arxiv.org/abs/1404.3873}{arXiv:1404.3873}.

\bibitem{OEIS}
The on-line encyclopedia of integer sequences, \url{https://oeis.org/}.

\bibitem{Pan}
Pan Y., Peelaers W., Exact {S}chur index in closed form,
 \href{https://doi.org/10.1103/physrevd.106.045017}{\textit{Phys. Rev.~D}}
 \textbf{106} (2022), 045017, 34~pages,
 \href{https://arxiv.org/abs/2112.09705}{arXiv:2112.09705}.

\bibitem{Qin}
Qin Z., Yu F., On {O}kounkov's conjecture connecting {H}ilbert schemes of
 points and multiple {$q$}-zeta values,
 \href{https://doi.org/10.1093/imrn/rnw244}{\textit{Int. Math. Res. Not.}}
 \textbf{2018} (2018), 321--361,
 \href{https://arxiv.org/abs/1510.00837}{arXiv:1510.00837}.

\bibitem{Edge}
Renteln P., The {H}ilbert series of the face ring of a flag complex,
 \href{https://doi.org/10.1007/s003730200045}{\textit{Graphs Combin.}}
 \textbf{18} (2002), 605--619.

\bibitem{Ittersum}
van Ittersum J.W.M., Partitions and quasimodular forms: {V}ariations on the
 {B}loch--{O}kounkov theorem, Ph.D. thesis, {U}trecht {U}niversity, 2021,
 \url{https://staff.fnwi.uva.nl/j.w.m.vanittersum/files/theses/thesis_vanIttersum_withcover.pdf}.

\bibitem{Zagier}
Zagier D., Multiple zeta values, {P}reprint, 1995.

\bibitem{Zhu}
Zhu Y., Modular invariance of characters of vertex operator algebras,
 \href{https://doi.org/10.1090/S0894-0347-96-00182-8}{\textit{J.~Amer. Math.
 Soc.}} \textbf{9} (1996), 237--302.

\bibitem{Zudilin2}
Zudilin W., Algebraic relations for multiple zeta values,
 \href{https://doi.org/10.1070/RM2003v058n01ABEH000592}{\textit{Russian Math.
 Surveys}} \textbf{58} (2003), 1--29.

\bibitem{Zudilin1}
Zudilin W., Multiple $q$-zeta brackets,
 \href{https://doi.org/10.3390/math3010119}{\textit{Mathematics}} \textbf{3}
 (2015), 119--130, \href{https://arxiv.org/abs/1412.0163}{arXiv:1412.0163}.

\end{thebibliography}
